\section[Notebook 1983-4: the cohomological equation on the torus]{Notebook 1983-4: the cohomological equation on the torus and Diophantine numbers\protect\verifie{Notebook1983_4}}\label{sec:quillen-1983-4}

\emph{Context.} In spring 1983, \carnet{1983}{4} studies the Connes algebra of the Kronecker foliation of irrational slope $\theta$ on the torus $T^2$ and its projective modules. On pages~71--72 (its pages~800--801), Quillen computes the Picard group of the topos of leaves of this foliation: the line bundles there are line bundles on $T^2$ equipped with a connection that is flat along the leaves, and the group sought is $H^1(T^2, \mathcal O^*)$, where $\mathcal O$ is the sheaf of smooth functions constant along the leaves.

\emph{The page.} The first step is $H^1(T^2, \mathcal O)$, computed by the fine resolution given by the operator $D = \theta\partial_x - \partial_y$ on $C^\infty(T^2)$. On the exponential $e^{2\pi i(mx + ny)}$, $D$ acts by the scalar $2\pi i(\theta m - n)$, which vanishes only at $(m, n) = (0, 0)$ since $\theta$ is irrational. Quillen deduces that $D$ is invertible off the constants, hence that $H^0$ and $H^1$ are both equal to $\mathbf C$, then, by the exponential sequence, that the Picard group is $\mathbf C^*/\{e^{2\pi i n\theta}\} \times \mathbf Z$.

\emph{Intuition.} Inverting $D$ on the exponentials is not enough: the series \[\sum b_{m,n}\,(2\pi i(\theta m - n))^{-1} e^{2\pi i(mx+ny)}\] must still define a smooth function. The coefficients of a smooth function decay faster than any power of $1 + |m| + |n|$; if the small divisors $|\theta m - n|$ are bounded below only by a power of $|m|^{-1}$, the division costs a fixed power and rapid decay survives. If on the contrary $\theta$ is very well approximated by rationals, one can place tiny coefficients $b_{m,n}$ exactly at the points where $\theta m - n$ is smaller still, and the quotient no longer decays. This is the phenomenon of small divisors, and it is the hypothesis that the page does not state.

\begin{enonce}
The sequences considered are indexed by $\mathbf Z^2$ and complex-valued. A sequence $a$ is \emph{rapidly decreasing} if, for every integer $k$, the quantity $|a_{m,n}|\,(1 + |m| + |n|)^k$ is bounded. A real number $\theta$ is \emph{Diophantine} if there exist $C > 0$ and an integer $N$ such that $|\theta m - n| \ge C\,|m|^{-N}$ for all integers $m \ne 0$ and $n$. The factor $2\pi$ is omitted: $D$ acts by $a_{m,n} \mapsto i(\theta m - n)\,a_{m,n}$.
\end{enonce}

\begin{encadre}
\begin{theoreme}[the edition's statement; missing hypothesis of notebook 1983-4, pp.~71--72]\label{thm:quillen-1983-4}
Let $\theta$ be an irrational real number. The operator $a \mapsto (i(\theta m - n)\,a_{m,n})_{m,n}$ sends rapidly decreasing sequences to rapidly decreasing sequences whose $(0,0)$ term is zero, and it is surjective onto the latter if and only if $\theta$ is Diophantine. More precisely: (1)~if $\theta$ is Diophantine (without any irrationality hypothesis) and $b$ is rapidly decreasing with $b_{0,0} = 0$, the sequence $a_{m,n} = b_{m,n}/(i(\theta m - n))$, $a_{0,0} = 0$, is rapidly decreasing and satisfies $i(\theta m - n)\,a_{m,n} = b_{m,n}$; (2)~if $\theta$ is irrational and not Diophantine, there exists a rapidly decreasing sequence $b$, with $b_{0,0} = 0$, that has no rapidly decreasing preimage.
In particular, the conclusion of the page holds for $\theta$ Diophantine, and fails for a Liouville number.
\end{theoreme}
\end{encadre}

\begin{proof}
Write $w_{m,n} = 1 + |m| + |n|$ and $\sigma_{m,n} = \theta m - n$. We have $|\sigma_{m,n}| \le (|\theta| + 1)\,w_{m,n}$, whence the first assertion.

(1) There exist $c > 0$ and $N$ such that $|\sigma_{m,n}|\,w_{m,n}^N \ge c$ for $(m, n) \ne (0,0)$: if $m \ne 0$, this is the hypothesis, since $|m| \le w_{m,n}$; if $m = 0$, we have $|\sigma_{0,n}| = |n| \ge 1$. Then $|a_{m,n}|\,w^k \le |b_{m,n}|\,w^{k+N}/c$, which is bounded. The lower bound involves only $|m|$, and $|m| \le w$: the dependence on $n$ is no obstacle in this direction.

(2) The negation of the hypothesis provides, for each $j$, a pair $q_j = (m_j, n_j)$ with $m_j \ne 0$ and $|\sigma_{q_j}| < (j+1)^{-1}|m_j|^{-2j}$; irrationality gives $\sigma_{q_j} \ne 0$. Put $b_{q} = \sigma_q / \sqrt{|\sigma_q|}$ if $q$ is one of the $q_j$, and $b_q = 0$ otherwise; thus $b_{0,0} = 0$ and $|b_{q_j}| = \sqrt{|\sigma_{q_j}|} < |m_j|^{-j}$. Here $n$ comes in: since $|\sigma_{q_j}| < 1$, we have $|n_j| \le |\theta|\,|m_j| + 1$, whence $w_{q_j} \le (3 + |\theta|)\,|m_j|$. For $j \ge k$, we get $|b_{q_j}|\,w_{q_j}^k \le (3 + |\theta|)^k |m_j|^{k-j} \le (3 + |\theta|)^k$; the indices $j < k$ are finitely many. Hence $b$ is rapidly decreasing. If $a$ were a preimage of it, we would have $|a_{q_j}| = |\sigma_{q_j}|^{-1/2} > \sqrt{j+1}$, and $a$ would not even be bounded.
\end{proof}

\begin{proposition}\label{prop:quillen-1983-4-liouville}
A Liouville number is not Diophantine. For such a $\theta$, there is therefore a rapidly decreasing sequence with zero $(0,0)$ term that is not in the image of $D$.
\end{proposition}

\begin{proof}
Let $C > 0$ and $N$. For infinitely many denominators $b$, there exists $a$ with $|\theta - a/b| < b^{-(N+2)}$; taking $b \ge 1/C$, we obtain $|\theta b - a| < b^{-(N+1)} \le C\,b^{-N}$, which contradicts the lower bound. A Liouville number is irrational, and Theorem~\ref{thm:quillen-1983-4} applies.
\end{proof}

\begin{remarque}
For $\theta$ non-Diophantine, the image of $D$ in the smooth functions of mean zero is dense but not closed: the quotient $H^1(T^2, \mathcal O)$ computed by the resolution is infinite-dimensional and non-Hausdorff. The computation of the Picard group that follows in the notebook therefore rests on the Diophantine hypothesis. The set of Diophantine reals has full measure, so that the conclusion of the page holds for almost every $\theta$.
\end{remarque}

\begin{sloppypar}
\noindent\emph{Verification in Lean.} The file \href{\leanblobmain/Notebook1983_4.lean}{\texttt{Notebook1983\_4.\allowbreak{}lean}} works at the level of sequences on $\mathbf Z \times \mathbf Z$, without Fourier analysis. It defines the weight \texttt{weight}, rapid decay \texttt{IsRapidDecay}, the predicate \texttt{IsDiophantine} and the symbol \texttt{symbol}. The forward direction is \texttt{IsDiophantine.isRapidDecay\_div} and \texttt{IsDiophantine.symbol\_mul\_div}, from the lower bound \texttt{IsDiophantine.lower\_bound}; the stability of rapid decay is \texttt{IsRapidDecay.symbol\_mul}. The converse is \texttt{exists\_not\_mem\_range\_of\_not\_isDiophantine}, the equivalence is \texttt{surjective\_iff\_isDiophantine}, and the case of Liouville numbers is \texttt{not\_isDiophantine\_of\_liouville} and \texttt{exists\_not\_mem\_range\_of\_liouville}, from mathlib's \texttt{Liouville} predicate. mathlib has neither rapidly decreasing sequences on $\mathbf Z^2$ nor Diophantine numbers in this sense; the file defines them.
\end{sloppypar}

Not formalized: the identification of smooth functions on $T^2$ with rapidly decreasing sequences of their Fourier coefficients, the resolution of the sheaf $\mathcal O$, the non-Hausdorff character of $H^1$ in the Liouville case, the full measure of the Diophantine numbers, and the computation of the Picard group.

\litterature{found in the literature: the phenomenon of small divisors and the cohomological equation on the torus are classical (V.~I.~Arnold, \emph{Geometrical Methods in the Theory of Ordinary Differential Equations}, supplementary chapters; M.~Herman, Publ. Math. IHÉS 49 (1979); A.~Katok, \emph{Combinatorial constructions in ergodic theory and dynamics}, AMS (2003)). For the foliated cohomology of the Kronecker foliation, which is $\mathbf C$ in degree~$1$ if and only if $\theta$ is Diophantine and infinite-dimensional non-Hausdorff otherwise: A.~El Kacimi Alaoui, Compositio Math. 49 (1983); A.~Haefliger, J. Differential Geom. 15 (1980). The statement is classical; the edition does not present it as new and only points out the hypothesis that the page does not state.}
